% Vietnamese translation for OI Public Library
% Source-File: IOI中国国家候选队论文/2007/day1/3.仇荣琦《欧拉回路性质与应用探究》.ppt
\documentclass[11pt]{article}
\usepackage{oipl-vn}

\newcommand{\Oh}{\mathcal{O}}
\newcommand{\indeg}{\operatorname{in}}
\newcommand{\outdeg}{\operatorname{out}}

\title{Bản trình chiếu: Tính chất và ứng dụng của chu trình Euler}
\author{Qiu Rongqi}
\date{2007}

\begin{document}
\maketitle

\origtitle{Slide companion for properties and applications of Euler circuits}
\sourcepath{IOI China candidate team papers / 2007 / day 1 / Qiu Rongqi.ppt}

\section{Phạm vi bản dịch}

Tệp PowerPoint đi kèm bài viết đã được dịch từ PDF về chu trình Euler. Các
slide dùng nhiều ví dụ bài tập, vì vậy bản ghi chú này gom lại điều kiện tồn
tại, thuật toán và các mô hình ứng dụng được nêu trong bài trình bày.

\section{Điều kiện tồn tại}

Với đồ thị vô hướng liên thông, chu trình Euler tồn tại khi và chỉ khi mọi
đỉnh đều có bậc chẵn. Đường đi Euler tồn tại khi và chỉ khi có đúng hai đỉnh
bậc lẻ; chúng là hai đầu của đường đi.

Với đồ thị có hướng, ta xét đồ thị nền để kiểm tra liên thông, đồng thời yêu
cầu \(\indeg(v)=\outdeg(v)\) với mọi đỉnh nếu cần chu trình Euler. Khi cần
đường đi Euler, hai đỉnh đặc biệt có chênh lệch bậc vào và bậc ra bằng \(1\)
theo hai chiều ngược nhau.

\section{Thuật toán}

Thuật toán chuẩn là Hierholzer: bắt đầu từ một đỉnh còn cạnh chưa dùng, đi
theo các cạnh chưa đi để tạo một chu trình, rồi chèn thêm các chu trình con
tại những đỉnh trên chu trình hiện tại còn cạnh chưa dùng. Nếu dùng danh sách
kề và xóa cạnh đã xét đúng cách, thời gian là \(\Oh(|E|)\).

\section{Ứng dụng}

Các slide nhắc nhiều bài mẫu như \textit{Depot Rearrangement},
\textit{Primitivus}, \textit{Ouroboros Snake}, \textit{Mosaic} và
\textit{Eclipse}. Điểm chung là bài toán ban đầu không nhất thiết hỏi trực
tiếp về chu trình Euler, nhưng sau khi mô hình hóa, yêu cầu trở thành đi qua
mỗi cạnh đúng một lần hoặc cân bằng lượng vào ra tại các đỉnh.

\section{Cách nhận diện mô hình}

Khi một bài toán yêu cầu sắp xếp các đối tượng sao cho mỗi quan hệ được dùng
đúng một lần, hoặc yêu cầu nối các mảnh theo tiền tố--hậu tố, chu trình Euler
là ứng viên tự nhiên. Với bài có chi phí, thường cần thêm bước biến đổi để
cân bằng bậc bằng các cạnh phụ hoặc đưa về bài toán đường đi ngắn nhất trước
khi dựng đường Euler cuối cùng.

\end{document}
